\originalchapter{来源与范围对应}
\originalsection{正文主线}
本表中的页码是公开旧讲义PDF页序, 与其印刷页码一致; 末页109为OCW条款. 分章PDF分别重复完整PDF相应正文, 另附条款页, 所以仅归档完整旧PDF以去重. 中文正文没有嵌入英文原页.
\begin{longtable}{p{.27\textwidth} p{.27\textwidth} p{.38\textwidth}}
\toprule 本书内容&旧讲义定位&编排与边界\\\midrule\endhead
第1章群与表示&§1.1--1.3、§1.7--1.8, 页5--14&群论必要语言与置换矩阵为衔接补充; 不展开一般代数分类\\
第2章完全可约&§1.3、§2.1、§3.1、§3.6, 页7--10、23、33--34、38--39&Schur、Maschke、平均内积与重数唯一性完整证明\\
第3章特征标&§2.6、§3.2--3.5, 页27--28、34--38&以Hom空间平均为统一证明; 置换轨道计数为补充\\
第4章正则与群代数&§2.2--2.5、§3.1、§3.7、§4.2, 页24--27、33、39--40、48--49&专门重证复群代数分块; 加入中心投影与反演; 不展开一般有限维代数根基\\
第5章经典表&§3.3、§3.8, 页35--36、40--42&循环群、$S_3$、$D_8$、$Q_8$、$A_4$, 补齐表的构造和完整性检查\\
第6章新表示&§1.10、§3.4、§3.9、§4.6--4.7, 页17--18、36--37、42--43、54&张量、对偶、对称与外平方、直积与虚表示\\
第7章诱导&§4.8--4.11, 页54--58&张量模型与函数模型互证; 两方向互反、传递性、投影公式\\
第8章双陪集&§4.9--4.11、§4.26, 页55--58、76--77&双陪集分解及判据为编者展开; 二面体群统一分类\\
第9章形式与$S_4$&§4.1、§3.8及§4.11, 页47--48、40--42、57--58&指标完整证明; $S_4$用商群与置换空间构造, 不调用Young图分类\\\bottomrule
\end{longtable}
前九章有编者例题与练习, 不冒充官方作业编号. 课程题解部分则以Homework和原题号保留身份, 共52道作业题及4道期末题. 旧讲义中的未布置练习并未全部收录, 官方未公开的题卷也不据书目或编号臆造.

\originalsection{课程题目与拓展范围}
Homework 1--4的题面来自旧PDF页10--15、18--22, 包含结合代数、Weyl代数、张量、Dehn不变量、$\mathfrak{sl}_2$及可解李代数. Homework 5--9来自页30--31、34、43--47、51、55、57--59、63、68, 包含扩张、群论应用、诱导、对称群及一般线性群局部题解. Homework 10--11来自页78--81、96--97, 涉及箭图与根系. 每题解答保留全部小问; 字符抽取损坏、原题笔误及需要补明的底域条件在对应位置说明.

期末作业单独官方PDF为2页, 第1页有4道题, 第2页为OCW条款. 本书分别处理路径代数、仿射群 $\GL_2(\mathbb F_q)\ltimes\mathbb F_q^2$、二元二十面体群的循环子群诱导, 以及 $\mathfrak{sl}_2$ 张量幂. 全部答案由编者推导并交叉复核; 官网无独立Solutions文件, 不能将本书答案称为官方标准答案.

课程题解为了处理原题而加入的局部理论, 不等于主源完整专题的系统覆盖. 李代数、箭图、根系、$\GL_2(\mathbb F_q)$和有限群半直积在题解中按需要展开; 范畴论、Gabriel定理全套证明、一般代数投射覆盖、完整Young图与Schur--Weyl理论、Artin定理、Frobenius整除性和Burnside可解性仍非本书完整专题.
